\subsection{Alpha Testing: Pruebas de Usabilidad y Control de Calidad con Usuarios Externos}

Durante el 10 de mayo de 2026, se realizó una prueba de usabilidad y control de calidad con usuarios externos, con el 
objetivo de validar la experiencia de usuario y la funcionalidad del juego. Para esto, se invitó a un total de 5 usuarios, 
que no formaban parte del equipo de desarrollo, a probar el juego y brindar su \textit{feedback} sobre la experiencia de 
usuario, la funcionalidad y la usabilidad del juego.
\vspace{0.25cm}

El proceso consistió una reunion en \textit{\textbf{Discord}}, en dicha reunion se les explico a los usuarios en que
consistia el juego, como utilizarlo y que funcionalidades tenia. Se les mostro el editor de mundos, para que pudieran crear 
su propio mundo. Luego se les dio acceso al editor de mundos, para que pudieran crear su propio mundo, y luego se les dio 
acceso al juego, para que pudieran jugar en el mundo que crearon.
\vspace{0.25cm}

Durante la prueba, se les pidió a los usuarios que brindaran su \textit{feedback} sobre la experiencia de usuario, la 
funcionalidad y la usabilidad del juego. A la vez, de que se brindaba soporte sobre dudas o dificultades que pudieran 
tener los usuarios durante la prueba.
\vspace{0.25cm}

Finalizada la prueba, se recopilo el \textit{feedback} de los usuarios, a traves de una encuesta en \textit{\textbf{Google Forms}}, 
donde se les preguntaba sobre la experiencia de usuario, la funcionalidad y la usabilidad del juego. Como asi tambien se les daba la 
posibilidad de brindar sugerencias o comentarios sobre el juego. Algunos de los comentarios, para mejorar, fueron:
\vspace{0.25cm}

\begin{itemize}
  \item \textit{"No se puede cambiar de pagina mientras cargan los assets"}
  \item \textit{"Falta de actividades / letras chicas"}
  \item \textit{"La UI es pésima y falta mucho feedback sobre qué herramienta estamos usando y qué caja de texto estamos llenando"}
\end{itemize}

Gracias al \textit{feedback} brindado por los usuarios, se pudieron identificar diferentes puntos de mejora en el juego, 
como la aspariencias de las \textit{UIs}, \textit{bugs} en la funcionalidad, y dificultades en la usabilidad del juego. 
El equipo tomo en cuenta el \textit{feedback} brindado por los usuarios, priorizando las mejoras a realizar, y se planifico 
el trabajo para implementar las mejoras identificadas en el juego.
\vspace{0.25cm}

En resumen, la prueba de usabilidad y control de calidad con usuarios externos, fue una experiencia muy positiva para el equipo, 
ya que permitió validar la experiencia de usuario y la funcionalidad del juego, y a la vez, identificar diferentes puntos de 
mejora en el juego, gracias al \textit{feedback} brindado por los usuarios.
